\documentclass[../Odyssey_Project.tex]{subfiles}

\begin{document}
\setcounter{section}{13}
\section{Characters}

\subsubsection*{Standing assumptions}
\begin{itemize}
    \item In this chapter $G$ denotes a finite group.
    \item All $\C G$-modules are finitely generated, equivalently, finite-dimensional as $\C$-vector spaces.
    \item $\C$ is the field of complex numbers.
\end{itemize}

\subsubsection*{Wedderburn structure of $\C G$}
\begin{itemize}
    \item Since $\C$ is algebraically closed of characteristic zero, Wedderburn theory pins down the structure of $\C G$ precisely.
    \item The simple summands $\mathcal{M}_{f_i}(\C)$ correspond bijectively to isomorphism classes of simple $\C G$-modules $S_i$, with $\dim_{\C} S_i = f_i$.
\end{itemize}

\begin{theorem}
\label{thm:14.1}
There is some $r \in \N$ and some $f_1, \ldots, f_r \in \N$ such that $\C G \cong \mathcal{M}_{f_1}(\C) \oplus \ldots \oplus \mathcal{M}_{f_r}(\C)$ as $\C$-algebras. Furthermore, there are exactly $r$ isomorphism classes of simple $\C G$-modules, and if we let $S_1, \ldots, S_r$ be representatives of these $r$ classes, then we can order the $S_i$ so that $\C G \cong f_1 S_1 \oplus \ldots \oplus f_r S_r$ as $\C G$-modules, where $\dim_{\C} S_i = f_i$ for each $i$. Any $\C G$-module can be written uniquely in the form $a_1 S_1 \oplus \ldots \oplus a_r S_r$, where the $a_i$ are non-negative integers.
\end{theorem}
\begin{proof}
    Since $\mathrm{char}(\C)=0$, Corollary~\ref{cor:12.8} shows that every finitely generated $\C G$-module is semisimple. In particular, the regular module $\C G$ is semisimple, and hence $\C G$ is a semisimple algebra by Lemma~\ref{lem:13.3}. Since $\C$ is algebraically closed, Theorem~\ref{thm:13.18} gives a $\C$-algebra isomorphism
    \[
        \C G \cong \mathcal{M}_{f_1}(\C) \oplus \cdots \oplus \mathcal{M}_{f_r}(\C)
    \]
    for some positive integers $f_1,\ldots,f_r$.\\
    For the algebra on the right, Theorem~\ref{thm:13.10} says that there are exactly $r$ isomorphism classes of simple modules, one coming from each matrix-algebra summand. More concretely, by Theorem~\ref{thm:13.6}, the simple module attached to the summand $\mathcal{M}_{f_i}(\C)$ is the natural column module $\C^{f_i}$, and the regular module $\mathcal{M}_{f_i}(\C)$ is the direct sum of $f_i$ copies of this module. Transporting these modules across the displayed algebra isomorphism gives simple $\C G$-modules $S_1,\ldots,S_r$ with $\dim_{\C}S_i=f_i$, and as $\C G$-modules we obtain
    \[
        \C G \cong f_1S_1 \oplus \cdots \oplus f_rS_r.
    \]
    Finally, every finitely generated $\C G$-module is semisimple, so it is a direct sum of simple modules. Since $S_1,\ldots,S_r$ represent all simple isomorphism classes, every such module has the form $a_1S_1 \oplus \cdots \oplus a_rS_r$ with $a_i \geq 0$. The multiplicities $a_i$ are uniquely determined by Proposition~\ref{prop:13.5}.
\end{proof}

\subsubsection*{Degrees of $G$}
\begin{itemize}
    \item The $\C$-dimensions $f_1, \ldots, f_r$ of the $r$ simple $\C G$-modules are called the \emph{degrees} of $G$.
    \item The trivial $\C G$-module $\C$ is one-dimensional and hence simple, so $G$ always has at least one degree equal to $1$; by convention we set $f_1 = 1$.
\end{itemize}

\begin{corollary}
\label{cor:14.2}
We have $\sum\limits_{i=1}^{r} f_i^2 = |G|$.
\end{corollary}
\begin{proof}
    Taking $\C$-dimensions in Theorem~\ref{thm:14.1} gives
    \[
        |G|=\dim_{\C}\C G
        =
        \dim_{\C}\left(\mathcal{M}_{f_1}(\C) \oplus \cdots \oplus \mathcal{M}_{f_r}(\C)\right)
        =
        \sum_{i=1}^{r}\dim_{\C}\mathcal{M}_{f_i}(\C)
        =
        \sum_{i=1}^{r}f_i^2.
    \]
\end{proof}

\subsubsection*{Degrees and conjugacy classes}
\begin{itemize}
    \item In fact the degrees of $G$ must divide $|G|$; we shall not use this fact, but we provide a proof of it in the Appendix.
    \item The next theorem links the number of simple $\C G$-modules to the conjugacy structure of $G$ via the center of $\C G$.
\end{itemize}

\begin{theorem}
\label{thm:14.3}
The number $r$ of simple $\C G$-modules is equal to the number of conjugacy classes of $G$.
\end{theorem}
\begin{proof}
    Let $Z$ be the center of the algebra $\C G$. By Theorem~\ref{thm:14.1}, $\C G$ is isomorphic with
    \[
        \mathcal{M}_{f_1}(\C) \oplus \cdots \oplus \mathcal{M}_{f_r}(\C).
    \]
    The center of a direct sum is the direct sum of the centers, and the center of $\mathcal{M}_{f_i}(\C)$ consists of the scalar matrices: commuting with the matrix units forces all off-diagonal entries to be zero and all diagonal entries to be equal. Hence $Z \cong \C^r$, so $\dim_{\C}Z=r$.\\
    We now compute the same dimension using the group basis of $\C G$. Write an arbitrary element of $\C G$ as
    \[
        a=\sum_{g \in G}\lambda_g g.
    \]
    If $a \in Z$, then $h a h^{-1}=a$ for every $h \in G$. Comparing the coefficients of the basis elements of $G$, we get
    \[
        \lambda_g=\lambda_{hgh^{-1}}
    \]
    for all $g,h \in G$. Thus the coefficients of a central element are constant on conjugacy classes. Conversely, if the coefficients are constant on conjugacy classes, then $h a h^{-1}=a$ for every $h \in G$, so $a$ commutes with every element of $G$, and hence with every element of $\C G$. Therefore $Z$ has as a basis the class sums
    \[
        \sum_{g \in K}g,
    \]
    where $K$ runs through the conjugacy classes of $G$. These class sums have disjoint supports in the basis $G$, so they are linearly independent. It follows that $\dim_{\C}Z$ is the number of conjugacy classes of $G$. Since we already proved that $\dim_{\C}Z=r$, the result follows.
\end{proof}

\subsubsection*{Defining characters}
\begin{itemize}
    \item In general there is no natural bijective correspondence between the conjugacy classes of $G$ and the simple $\C G$-modules. However, if $G$ is a symmetric group, then there is such a correspondence, although we shall not develop it here; for an overview, see [13, Section 4.1].
    \item If $U$ is a $\C G$-module, then each $g \in G$ defines an invertible linear transformation of $U$ that sends $u \in U$ to $gu$. The \emph{character} of $U$ is the function $\chi_U \colon G \to \C$, where $\chi_U(g)$ is the trace of this linear transformation of $U$ defined by $g$.
    \item For example, $\chi_U(1) = \dim_{\C} U$, since the identity element of $G$ induces the identity transformation on $U$.
    \item If $\rho \colon G \to \mathrm{GL}(U)$ is the representation corresponding to $U$, then $\chi_U(g)$ is just the trace of the map $\rho(g)$.
    \item Isomorphic $\C G$-modules have equal characters.
    \item For any $g, h \in G$ the linear transformations of $U$ defined by $g$ and $hgh^{-1}$ are similar and hence have the same trace; so any character of $G$ is constant on each conjugacy class of $G$.
    \item For example, let $U = \C G$ and let $g \in G$. With respect to the basis $G$ of $\C G$, $\chi_U(g)$ equals the number of elements $x \in G$ for which $gx = x$. Therefore $\chi_U(1) = |G|$ and $\chi_U(g) = 0$ for every $1 \neq g \in G$. This character is called the \emph{regular character} of $G$.
    \item The theory of characters was developed by Frobenius and others, starting in 1896. At first nothing was known about linear representations, let alone modules over the group algebra; Frobenius defined characters as being functions from $G$ to $\C$ satisfying certain properties, but it turned out that his characters were exactly the trace functions of finitely generated $\C G$-modules.
    \item We will denote the characters of the $r$ simple $\C G$-modules by $\chi_1, \ldots, \chi_r$; these characters will be referred to as the \emph{irreducible characters} of $G$. Whenever we say that $S_1, \ldots, S_r$ are the distinct simple $\C G$-modules, we implicitly order them so that $\chi_{S_i} = \chi_i$ for each $i$.
    \item In keeping with the convention that $f_1$ denotes the degree of the trivial representation, we let $\chi_1$ be the character of the trivial representation; we call $\chi_1$ the \emph{principal character} of $G$, and we have $\chi_1(g) = 1$ for all $g \in G$.
    \item A character of a one-dimensional $\C G$-module is called a \emph{linear character}. Since one-dimensional modules are simple, all linear characters are irreducible. If $\chi$ arises from a one-dimensional $\C G$-module $U$, then for $g, h \in G$ and $u \in U$ we have $gu = \chi(g)u$ and $hu = \chi(h)u$, so $\chi(gh)u = (gh)u = \chi(g)\chi(h)u$; hence $\chi$ is a homomorphism from $G$ to $\C^{\times}$. Conversely, given a homomorphism $\varphi \colon G \to \C^{\times}$, we can define a one-dimensional $\C G$-module $U$ by $gu = \varphi(g) u$ for $g \in G$ and $u \in U$, and we then have $\chi_U = \varphi$. Therefore, linear characters of $G$ are exactly the same as group homomorphisms from $G$ to $\C^{\times}$.
\end{itemize}

\begin{proposition}
\label{prop:14.4}
Let $U$ be a $\C G$-module, let $\rho \colon G \to \mathrm{GL}(U)$ be the representation corresponding to $U$, and let $g \in G$ be of order $n$. Then:
\begin{enumerate}
    \item[(i)] $\rho(g)$ is diagonalizable.
    \item[(ii)] $\chi_U(g)$ equals the sum (with multiplicities) of the eigenvalues of $\rho(g)$.
    \item[(iii)] $\chi_U(g)$ is a sum of $\chi_U(1)$ $n$th roots of unity.
    \item[(iv)] $\chi_U(g^{-1}) = \overline{\chi_U(g)}$. (Here $\overline{z}$ denotes the complex conjugate of $z \in \C$.)
    \item[(v)] $|\chi_U(g)| \leq \chi_U(1)$.
    \item[(vi)] $\{x \in G \mid \chi_U(x) = \chi_U(1)\}$ is a normal subgroup of $G$.
\end{enumerate}
\end{proposition}
\begin{proof}
    Since $g^n=1$, we have $\rho(g)^n=I$, so the minimal polynomial of $\rho(g)$ divides $x^n-1$. Over $\C$, the polynomial $x^n-1$ splits into distinct linear factors. Hence the minimal polynomial of $\rho(g)$ also splits into distinct linear factors, and therefore $\rho(g)$ is diagonalizable. This proves (i).\\
    Once $\rho(g)$ is diagonalizable, we may choose a basis of $U$ consisting of eigenvectors for $\rho(g)$. In such a basis the matrix of $\rho(g)$ is diagonal, so its trace is the sum of its eigenvalues, counted with multiplicity. This proves (ii). Since every eigenvalue $\lambda$ satisfies $\lambda^n=1$, each eigenvalue is an $n$th root of unity. There are $\dim_{\C}U=\chi_U(1)$ eigenvalues with multiplicity, so $\chi_U(g)$ is a sum of $\chi_U(1)$ $n$th roots of unity. This proves (iii).\\
    The eigenvalues of $\rho(g^{-1})=\rho(g)^{-1}$ are the inverses of the eigenvalues of $\rho(g)$. If $\lambda$ is a root of unity, then $\lambda^{-1}=\overline{\lambda}$. Hence the eigenvalues of $\rho(g^{-1})$ are the complex conjugates of the eigenvalues of $\rho(g)$, and taking traces gives $\chi_U(g^{-1})=\overline{\chi_U(g)}$. This proves (iv).\\
    Part (iii) writes $\chi_U(g)=\lambda_1+\cdots+\lambda_m$, where $m=\chi_U(1)$ and each $|\lambda_i|=1$. Therefore
    \[
        |\chi_U(g)| \leq |\lambda_1|+\cdots+|\lambda_m|=m=\chi_U(1),
    \]
    proving (v).\\
    It remains to prove (vi). Let
    \[
        K=\{x \in G \mid \chi_U(x)=\chi_U(1)\}.
    \]
    For $x \in G$, the preceding parts write $\chi_U(x)$ as a sum of $\chi_U(1)$ roots of unity. Such a sum can equal the positive real number $\chi_U(1)$ only when every summand is $1$. Since $\rho(x)$ is diagonalizable, this happens exactly when $\rho(x)=I$. Thus $K=\ker\rho$, and hence $K \trianglelefteq G$.
\end{proof}

\subsubsection*{Algebra of characters}
\begin{itemize}
    \item If $\chi$ and $\psi$ are characters of $G$, we can define new functions $\chi + \psi$ and $\chi\psi$ from $G$ to $\C$ by $(\chi + \psi)(g) = \chi(g) + \psi(g)$ and $(\chi\psi)(g) = \chi(g)\psi(g)$ for $g \in G$.
    \item These new functions obtained from characters are not, a priori, characters themselves.
    \item Given a scalar $\lambda \in \C$, we can also define a new function $\lambda \chi \colon G \to \C$ by $(\lambda \chi)(g) = \lambda \chi(g)$, so we view characters of $G$ as elements of a $\C$-vector space of functions from $G$ to $\C$.
\end{itemize}

\begin{proposition}
\label{prop:14.5}
The irreducible characters of $G$ are, as functions from $G$ to $\C$, linearly independent over $\C$.
\end{proposition}
\begin{proof}
    For each irreducible character $\chi_i$, extend $\chi_i$ linearly from the basis $G$ of $\C G$ to a linear map $\C G \to \C$. Equivalently, $\chi_i(a)$ is the trace of the linear transformation of $S_i$ defined by $a \in \C G$.\\
    Using the algebra decomposition from Theorem~\ref{thm:14.1}, let $e_i$ be the element that is the identity matrix in the $i$th summand $\mathcal{M}_{f_i}(\C)$ and zero in all other summands. Then $e_i$ acts as the identity on $S_i$, so $\chi_i(e_i)=\dim_{\C}S_i=f_i$. If $j \neq i$, then $e_i$ acts as zero on $S_j$, so $\chi_j(e_i)=0$.\\
    Suppose that $\lambda_1,\ldots,\lambda_r \in \C$ satisfy
    \[
        \lambda_1\chi_1+\cdots+\lambda_r\chi_r=0
    \]
    as functions on $G$. Since $G$ is a basis of $\C G$, the same equality holds after extending the characters linearly to $\C G$. Evaluating at $e_i$ gives
    \[
        0=\sum_{j=1}^{r}\lambda_j\chi_j(e_i)=\lambda_i f_i.
    \]
    Since $f_i>0$, we get $\lambda_i=0$. This holds for every $i$, so the irreducible characters are linearly independent.
\end{proof}

\begin{lemma}
\label{lem:14.6}
$\chi_{U \oplus V} = \chi_U + \chi_V$ for any $\C G$-modules $U$ and $V$.
\end{lemma}
\begin{proof}
    Fix $g \in G$. Choose a basis of $U \oplus V$ whose first vectors form a basis of $U \oplus 0$ and whose remaining vectors form a basis of $0 \oplus V$. With respect to this basis, the linear transformation induced by $g$ has block diagonal form
    \[
        \begin{pmatrix}
            \rho_U(g) & 0\\
            0 & \rho_V(g)
        \end{pmatrix}.
    \]
    The trace of a block diagonal matrix is the sum of the traces of its diagonal blocks. Therefore
    \[
        \chi_{U \oplus V}(g)=\chi_U(g)+\chi_V(g).
    \]
    Since this holds for every $g \in G$, we have $\chi_{U \oplus V}=\chi_U+\chi_V$.
\end{proof}

\subsubsection*{Characters distinguish modules}
\begin{itemize}
    \item Characters of $\C G$-modules suffice to distinguish $\C G$-modules up to isomorphism.
\end{itemize}

\begin{theorem}
\label{thm:14.7}
If $S_1, \ldots, S_r$ are the distinct simple $\C G$-modules, then the character of the $\C G$-module $a_1 S_1 \oplus \ldots \oplus a_r S_r$ (where the $a_i$ are non-negative integers) is $a_1 \chi_1 + \ldots + a_r \chi_r$. Consequently, two $\C G$-modules are isomorphic iff their characters are equal.
\end{theorem}
\begin{proof}
    Repeatedly applying Lemma~\ref{lem:14.6} gives
    \[
        \chi_{a_1S_1 \oplus \cdots \oplus a_rS_r}
        =
        a_1\chi_1+\cdots+a_r\chi_r.
    \]
    Thus the displayed formula for the character of a decomposed module holds.\\
    Isomorphic modules have equal characters, so only the converse needs proof. Suppose that $U$ and $V$ are $\C G$-modules with $\chi_U=\chi_V$. By Theorem~\ref{thm:14.1}, write
    \[
        U \cong a_1S_1 \oplus \cdots \oplus a_rS_r,
        \qquad
        V \cong b_1S_1 \oplus \cdots \oplus b_rS_r.
    \]
    Taking characters and using the first part of the proof gives
    \[
        0=\chi_U-\chi_V=(a_1-b_1)\chi_1+\cdots+(a_r-b_r)\chi_r.
    \]
    Proposition~\ref{prop:14.5} says that the $\chi_i$ are linearly independent, so $a_i-b_i=0$ for every $i$. Hence $a_i=b_i$ for every $i$, and therefore $U \cong V$.
\end{proof}

\subsubsection*{Tensor products, duals, and Hom}
\begin{itemize}
    \item Although each character determines a $\C G$-module up to isomorphism by Theorem~\ref{thm:14.7}, there is no generic way to construct a module from its corresponding character, and so in some sense information is lost by studying characters in lieu of modules.
    \item Characters turn out to be an efficient means of translating information about the ordinary representation theory of $G$ to information about $G$ itself; from now on we focus on characters rather than modules.
    \item Recall from Section~12 that if $U$ and $V$ are $\C G$-modules, then $U \otimes_{\C} V$ and $\mathrm{Hom}_{\C}(U,V)$, written here as $U \otimes V$ and $\mathrm{Hom}(U,V)$, admit natural $\C G$-module structures.
\end{itemize}

\begin{proposition}
\label{prop:14.8}
Let $U$ and $V$ be $\C G$-modules. Then:
\begin{enumerate}
    \item[(i)] $\chi_{U \otimes V} = \chi_U \chi_V$.
    \item[(ii)] $\chi_{U^*} = \overline{\chi_U}$. (Recall that $U^* = \mathrm{Hom}_{\C}(U, \C)$.)
    \item[(iii)] $\chi_{\mathrm{Hom}(U,V)} = \overline{\chi_U} \chi_V$.
\end{enumerate}
\end{proposition}
\begin{proof}
    We prove the three assertions separately.\\
    (i) Fix $g \in G$. By Proposition~\ref{prop:14.4}, the transformations induced by $g$ on $U$ and $V$ are diagonalizable. Choose eigenbases $u_1,\ldots,u_m$ of $U$ and $v_1,\ldots,v_n$ of $V$, with
    \[
        gu_i=\lambda_i u_i,
        \qquad
        gv_j=\mu_j v_j.
    \]
    Then $u_i \otimes v_j$, with $1 \leq i \leq m$ and $1 \leq j \leq n$, is a basis of $U \otimes V$, and
    \[
        g(u_i \otimes v_j)=gu_i \otimes gv_j
        =
        \lambda_i\mu_j(u_i \otimes v_j).
    \]
    Hence the eigenvalues of $g$ on $U \otimes V$ are the products $\lambda_i\mu_j$, and so
    \[
        \chi_{U \otimes V}(g)
        =
        \sum_{i,j}\lambda_i\mu_j
        =
        \left(\sum_i\lambda_i\right)\left(\sum_j\mu_j\right)
        =
        \chi_U(g)\chi_V(g).
    \]
    This proves (i).\\
    (ii) Again fix $g \in G$, and choose an eigenbasis $u_1,\ldots,u_m$ of $U$ with $gu_i=\lambda_i u_i$. Let $\varphi_1,\ldots,\varphi_m$ be the dual basis of $U^*$. Since each $\lambda_i$ is a root of unity, $\lambda_i^{-1}=\overline{\lambda_i}$. For each $i,j$ we have
    \[
        (g\varphi_i)(u_j)
        =
        \varphi_i(g^{-1}u_j)
        =
        \varphi_i(\lambda_j^{-1}u_j)
        =
        \overline{\lambda_j}\delta_{ij}.
    \]
    Thus $g\varphi_i=\overline{\lambda_i}\varphi_i$. Therefore the eigenvalues of $g$ on $U^*$ are $\overline{\lambda_1},\ldots,\overline{\lambda_m}$, and
    \[
        \chi_{U^*}(g)=\overline{\lambda_1}+\cdots+\overline{\lambda_m}
        =
        \overline{\chi_U(g)}.
    \]
    This proves (ii).\\
    (iii) Proposition~\ref{prop:12.7} gives an isomorphism of $\C G$-modules
    \[
        \mathrm{Hom}(U,V) \cong U^* \otimes V.
    \]
    Isomorphic modules have equal characters, so by (i) and (ii),
    \[
        \chi_{\mathrm{Hom}(U,V)}
        =
        \chi_{U^* \otimes V}
        =
        \chi_{U^*}\chi_V
        =
        \overline{\chi_U}\chi_V.
    \]
\end{proof}

\subsubsection*{Virtual characters and class functions}
\begin{itemize}
    \item A \emph{virtual character} of $G$ is a $\Z$-linear combination of the irreducible characters of $G$. (Some authors prefer the term ``generalized character.'') Characters are virtual characters by Theorem~\ref{thm:14.7}.
\end{itemize}

\begin{corollary}
\label{cor:14.9}
The virtual characters of $G$ form a ring.
\end{corollary}
\begin{proof}
    Addition and additive inverses are immediate from the definition of virtual characters as $\Z$-linear combinations of the irreducible characters. The only point to check is closure under multiplication.\\
    If $\chi$ and $\psi$ are ordinary characters, then Proposition~\ref{prop:14.8} shows that $\chi\psi$ is the character of a tensor product module. Hence $\chi\psi$ is again an ordinary character, and by Theorem~\ref{thm:14.7} it is a non-negative integer combination of the irreducible characters. Now if
    \[
        \alpha=\sum_i a_i\chi_i,
        \qquad
        \beta=\sum_j b_j\chi_j
    \]
    are virtual characters, then
    \[
        \alpha\beta=\sum_{i,j}a_i b_j(\chi_i\chi_j),
    \]
    and each product $\chi_i\chi_j$ is an integer combination of irreducible characters by the preceding paragraph. Thus $\alpha\beta$ is again a virtual character. The principal character is the multiplicative identity, so the virtual characters form a ring.
\end{proof}

\subsubsection*{Class functions}
\begin{itemize}
    \item A \emph{class function} on $G$ is a function from $G$ to $\C$ whose value within any conjugacy class is constant.
    \item Characters of $\C G$-modules are class functions.
    \item The set of class functions on $G$ forms a $\C$-vector space of dimension $r$, where $r$ is the number of conjugacy classes of $G$; an obvious basis is the set of functions taking the value $1$ on a single conjugacy class and $0$ on all other classes.
\end{itemize}

\begin{proposition}
\label{prop:14.10}
The irreducible characters of $G$ form a basis for the space of class functions on $G$.
\end{proposition}
\begin{proof}
    Every irreducible character is a character, and hence is a class function. By Proposition~\ref{prop:14.5}, the irreducible characters $\chi_1,\ldots,\chi_r$ are linearly independent. By Theorem~\ref{thm:14.3}, the number $r$ is also the number of conjugacy classes of $G$, which is the dimension of the space of class functions on $G$. Therefore we have $r$ linearly independent vectors in an $r$-dimensional vector space, so the irreducible characters form a basis.
\end{proof}

\subsubsection*{Inner product on class functions}
\begin{itemize}
    \item If $\alpha$ and $\beta$ are class functions on $G$, then their \emph{inner product} is the complex number
    \[
        (\alpha, \beta) = \frac{1}{|G|} \sum_{g \in G} \alpha(g) \overline{\beta(g)}.
    \]
    \item This function $(\,\cdot\,, \,\cdot\,)$ is, as the name suggests, an inner product on the space of class functions:
    \begin{itemize}
        \item $(\alpha, \alpha) \geq 0$ for all $\alpha$, and $(\alpha, \alpha) = 0$ iff $\alpha = 0$.
        \item $(\alpha, \beta) = \overline{(\beta, \alpha)}$ for all $\alpha, \beta$.
        \item $(\lambda \alpha, \beta) = \lambda(\alpha, \beta)$ for all $\alpha, \beta$ and all $\lambda \in \C$.
        \item $(\alpha_1 + \alpha_2, \beta) = (\alpha_1, \beta) + (\alpha_2, \beta)$ for all $\alpha_1, \alpha_2, \beta$.
    \end{itemize}
    \item The following are easy consequences of the above properties:
    \begin{itemize}
        \item $(\alpha, \lambda \beta) = \overline{\lambda}(\alpha, \beta)$ for all $\alpha, \beta$ and all $\lambda \in \C$.
        \item $(\alpha, \beta_1 + \beta_2) = (\alpha, \beta_1) + (\alpha, \beta_2)$ for all $\alpha, \beta_1, \beta_2$.
    \end{itemize}
    \item We need a small piece of notation: if $U$ is a $\C G$-module, the set of elements of $U$ on which $G$ acts trivially is a $\C G$-submodule, denoted $U^G = \{u \in U \mid gu = u \text{ for all } g \in G\}$.
\end{itemize}

\begin{lemma}
\label{lem:14.11}
If $U$ is a $\C G$-module, then $\dim_{\C} U^G = \frac{1}{|G|} \sum\limits_{g \in G} \chi_U(g)$.
\end{lemma}
\begin{proof}
    Let
    \[
        a=\frac{1}{|G|}\sum_{g \in G}g \in \C G,
    \]
    and let $T$ be the linear transformation of $U$ defined by the action of $a$. For any $h \in G$, we have $ha=a$, because left multiplication by $h$ merely permutes the elements of $G$ inside the sum. Also,
    \[
        a^2
        =
        \frac{1}{|G|^2}\sum_{g,k \in G}gk
        =
        \frac{1}{|G|}\sum_{y \in G}y
        =
        a,
    \]
    since each $y \in G$ occurs as $gk$ exactly $|G|$ times. Hence $T^2=T$. The polynomial $x^2-x=x(x-1)$ has distinct roots, so $T$ is diagonalizable with eigenvalues only $0$ and $1$.\\
    Let $U_1$ be the eigenspace of $T$ with eigenvalue $1$. If $u \in U_1$, then for every $h \in G$ we have
    \[
        hu=hTu=(ha)u=au=Tu=u,
    \]
    so $u \in U^G$. Conversely, if $u \in U^G$, then
    \[
        Tu=au=\frac{1}{|G|}\sum_{g \in G}gu
        =
        \frac{1}{|G|}\sum_{g \in G}u
        =
        u,
    \]
    so $u \in U_1$. Thus $U^G=U_1$.\\
    Since $T$ is diagonalizable with eigenvalues $0$ and $1$, its trace is the dimension of its $1$-eigenspace, namely $\dim_{\C}U^G$. On the other hand, by linearity of trace,
    \[
        \mathrm{tr}(T)
        =
        \frac{1}{|G|}\sum_{g \in G}\chi_U(g).
    \]
    Therefore
    \[
        \dim_{\C}U^G=\frac{1}{|G|}\sum_{g \in G}\chi_U(g).
    \]
\end{proof}

\begin{theorem}
\label{thm:14.12}
We have $(\chi_U, \chi_V) = \dim_{\C} \mathrm{Hom}_{\C G}(U, V)$ for any $\C G$-modules $U$ and $V$.
\end{theorem}
\begin{proof}
    First we identify a fixed-point space. Inside the $\C G$-module $\mathrm{Hom}(U,V)$, the $\C G$-homomorphisms are exactly the $G$-fixed maps. Recall that the $G$-action is given by
    \[
        (g\varphi)(u)=g\varphi(g^{-1}u).
    \]
    If $\varphi \in \mathrm{Hom}_{\C G}(U,V)$, then
    \[
        (g\varphi)(u)=g\varphi(g^{-1}u)=gg^{-1}\varphi(u)=\varphi(u),
    \]
    so $\varphi$ is fixed by every $g \in G$. Conversely, if $g\varphi=\varphi$ for every $g \in G$, then applying this equality to $gu$ gives
    \[
        \varphi(gu)=(g\varphi)(gu)=g\varphi(u),
    \]
    so $\varphi$ is $G$-linear. Hence
    \[
        \mathrm{Hom}_{\C G}(U,V)=\mathrm{Hom}(U,V)^G.
    \]
    By Lemma~\ref{lem:14.11} and Proposition~\ref{prop:14.8}, we now get
    \[
        \dim_{\C}\mathrm{Hom}_{\C G}(U,V)
        =
        \frac{1}{|G|}\sum_{g \in G}\chi_{\mathrm{Hom}(U,V)}(g)
        =
        \frac{1}{|G|}\sum_{g \in G}\overline{\chi_U(g)}\chi_V(g)
        =
        (\chi_V,\chi_U).
    \]
    The right-hand side is equal to a dimension, so it is a real number. By conjugate symmetry of the inner product,
    \[
        (\chi_U,\chi_V)=\overline{(\chi_V,\chi_U)}=(\chi_V,\chi_U).
    \]
    Therefore $(\chi_U,\chi_V)=\dim_{\C}\mathrm{Hom}_{\C G}(U,V)$.
\end{proof}

\subsection*{Exercises}

Throughout these exercises, $G$ denotes a finite group, and all modules are finitely generated.

\begin{exercise}
\label{ex:14.1}
Let $A$ be a semisimple $\C$-algebra. Let $[A, A]$ be the subspace of $A$ spanned by the set $\{ab - ba \mid a, b \in A\}$. Show that the codimension of $[A, A]$ in $A$ (which is just the dimension of $A/[A, A]$) is equal to the number of isomorphism classes of simple $A$-modules.
\end{exercise}
\begin{solution}
    By Theorem~\ref{thm:13.18}, we may write
    \[
        A\cong \mathcal{M}_{n_1}(\C)\oplus\cdots\oplus \mathcal{M}_{n_r}(\C)
    \]
    as $\C$-algebras. By Theorem~\ref{thm:13.10}, this algebra has exactly $r$ isomorphism classes of simple modules. Thus it remains to show that the quotient by the commutator subspace has dimension $r$.\\
    First consider $B=\mathcal{M}_n(\C)$. Every commutator in $B$ has trace zero, since
    \[
        \operatorname{tr}(XY-YX)=\operatorname{tr}(XY)-\operatorname{tr}(YX)=0.
    \]
    Conversely, the trace-zero matrices are spanned by commutators. Indeed, for $i\neq j$,
    \[
        E_{ij}=E_{ii}E_{ij}-E_{ij}E_{ii},
    \]
    and for $i\neq j$,
    \[
        E_{ii}-E_{jj}=E_{ij}E_{ji}-E_{ji}E_{ij}.
    \]
    The off-diagonal matrix units together with the diagonal differences span the trace-zero subspace of $\mathcal{M}_n(\C)$. Hence $[B,B]$ is exactly the trace-zero subspace, and so $B/[B,B]$ is one-dimensional.\\
    In a direct sum of algebras, commutators are computed componentwise, so
    \[
        [A,A]\cong [\mathcal{M}_{n_1}(\C),\mathcal{M}_{n_1}(\C)]\oplus\cdots\oplus [\mathcal{M}_{n_r}(\C),\mathcal{M}_{n_r}(\C)].
    \]
    Therefore each matrix-algebra summand contributes one dimension to the quotient, and
    \[
        \dim_{\C}A/[A,A]=r.
    \]
    This is exactly the number of isomorphism classes of simple $A$-modules.
\end{solution}

\begin{exercise}
\label{ex:14.2}
(cont.) Let $A = \C G$, and let $g_1, \ldots, g_r$ be representatives of the $r$ conjugacy classes of $G$. Show that $\{g_i + [A, A] \mid 1 \leq i \leq r\}$ is a basis of $A/[A, A]$.
\end{exercise}
\begin{solution}
    Let $\pi \colon A\to A/[A,A]$ be the quotient map. If $g,h\in G$, then
    \[
        hgh^{-1}-g=(hg)h^{-1}-h^{-1}(hg)\in [A,A].
    \]
    Thus conjugate elements of $G$ have the same image in $A/[A,A]$. Since every element of $A=\C G$ is a $\C$-linear combination of elements of $G$, it follows that the cosets
    \[
        g_i+[A,A],\qquad 1\leq i\leq r,
    \]
    span $A/[A,A]$.\\
    By Exercise~\ref{ex:14.1}, the dimension of $A/[A,A]$ is the number of isomorphism classes of simple $\C G$-modules. By Theorem~\ref{thm:14.3}, this number is the number of conjugacy classes of $G$, namely $r$. Thus we have $r$ spanning vectors in an $r$-dimensional vector space, and hence
    \[
        \{g_i+[A,A]\mid 1\leq i\leq r\}
    \]
    is a basis of $A/[A,A]$.
\end{solution}

\begin{exercise}
\label{ex:14.3}
Show that if $S$ is a simple $\C G$-module and $U$ is a one-dimensional $\C G$-module, then $S \otimes U$ is simple.
\end{exercise}
\begin{solution}
    Choose $0\neq u\in U$. Since $U$ is one-dimensional, for each $g\in G$ there is a scalar $\lambda_g\in \C^{\times}$ such that $gu=\lambda_g u$. Let $0\neq W\leq S\otimes U$ be a $\C G$-submodule. Under the vector-space isomorphism
    \[
        S\to S\otimes U,\qquad s\mapsto s\otimes u,
    \]
    let $T$ be the inverse image of $W$. Thus
    \[
        T=\{s\in S\mid s\otimes u\in W\}.
    \]
    Since $W\neq0$, we have $T\neq0$. If $s\in T$ and $g\in G$, then
    \[
        g(s\otimes u)=gs\otimes gu=gs\otimes \lambda_g u=\lambda_g(gs\otimes u)
    \]
    lies in $W$. Since $\lambda_g\neq0$, it follows that $gs\otimes u\in W$, so $gs\in T$. Hence $T$ is a non-zero $\C G$-submodule of $S$. Since $S$ is simple, $T=S$. Therefore $W=S\otimes U$, and so $S\otimes U$ is simple.
\end{solution}

\begin{exercise}
\label{ex:14.4}
Show that the set of isomorphism classes of the one-dimensional $\C G$-modules forms a group under the taking of tensor products.
\end{exercise}
\begin{solution}
    For one-dimensional $\C G$-modules $U$ and $V$, define
    \[
        [U][V]=[U\otimes V],
    \]
    where brackets denote isomorphism classes. This is well-defined because isomorphic modules give isomorphic tensor products, and $U\otimes V$ is again one-dimensional. Associativity follows from the usual associativity isomorphism for tensor products.\\
    The identity element is the class of the trivial one-dimensional module $\C$, since $U\otimes \C\cong U$ as $\C G$-modules. The inverse of $[U]$ is $[U^*]$. Indeed, the evaluation map
    \[
        U\otimes U^*\to \C,\qquad u\otimes \varphi\mapsto \varphi(u),
    \]
    is a $\C G$-module homomorphism, because
    \[
        (g\varphi)(gu)=\varphi(g^{-1}gu)=\varphi(u).
    \]
    Both $U\otimes U^*$ and $\C$ are one-dimensional, and the evaluation map is non-zero, so it is an isomorphism. Thus every element has an inverse, and the isomorphism classes of one-dimensional $\C G$-modules form a group under tensor product.
\end{solution}

\begin{exercise}
\label{ex:14.5}
Let $\chi$ be an irreducible character of $G$. Show that if $\lambda$ is any $|G|$th root of unity, then $\{x \in G \mid \chi(x) = \lambda \chi(1)\} \trianglelefteq G$.
\end{exercise}
\begin{solution}
    As written, this assertion is false. Let $G=\langle a\rangle$ be cyclic of order $2$, and let $U$ be the one-dimensional $\C G$-module on which $a$ acts as multiplication by $-1$. Its character $\chi$ is irreducible, with
    \[
        \chi(1)=1,\qquad \chi(a)=-1.
    \]
    Taking $\lambda=-1$, which is a $|G|$th root of unity, gives
    \[
        \{x\in G\mid \chi(x)=\lambda\chi(1)\}=\{a\}.
    \]
    This set does not contain the identity element, and hence is not a subgroup of $G$.\\
    The nearby true statement is that this set is stable under conjugation. Indeed, if $\rho$ affords $\chi$ and $\chi(x)=\lambda\chi(1)$, then the eigenvalues of $\rho(x)$ are roots of unity and have absolute value $1$. Their sum has absolute value $\chi(1)$, so equality in the triangle inequality forces all of them to be $\lambda$. Hence $\rho(x)=\lambda I$, and conjugating $x$ inside $G$ leaves this scalar matrix unchanged. For $\lambda=1$, the set is $\ker\rho$, which is the normal subgroup from Proposition~\ref{prop:14.4}.
\end{solution}

\begin{exercise}
\label{ex:14.6}
Show that a $\C G$-module $U$ is simple iff its dual $U^*$ is simple. Conclude that the complex conjugate of an irreducible character is again an irreducible character.
\end{exercise}
\begin{solution}
    Suppose first that $U$ is simple. Let $0\neq W\leq U^*$ be a $\C G$-submodule, and define
    \[
        W^{\perp}=\{u\in U\mid \varphi(u)=0\text{ for all }\varphi\in W\}.
    \]
    We claim that $W^{\perp}$ is a $\C G$-submodule of $U$. If $u\in W^{\perp}$, $g\in G$, and $\varphi\in W$, then
    \[
        \varphi(gu)=(g^{-1}\varphi)(u)=0,
    \]
    because $g^{-1}\varphi\in W$. Hence $gu\in W^{\perp}$. Since $W\neq0$, some $\varphi\in W$ is non-zero, so $W^{\perp}\neq U$. By simplicity of $U$, we get $W^{\perp}=0$. Since $U$ is finite-dimensional,
    \[
        \dim_{\C}W^{\perp}=\dim_{\C}U-\dim_{\C}W,
    \]
    and therefore $\dim_{\C}W=\dim_{\C}U=\dim_{\C}U^*$. Thus $W=U^*$, and $U^*$ is simple.\\
    Conversely, suppose that $U^*$ is simple. Applying the first direction to $U^*$ shows that $U^{**}$ is simple. The natural map
    \[
        U\to U^{**},\qquad u\mapsto (\varphi\mapsto \varphi(u)),
    \]
    is a $\C G$-module isomorphism, since
    \[
        (g(\varphi\mapsto \varphi(u)))(\varphi)=(g^{-1}\varphi)(u)=\varphi(gu).
    \]
    Hence $U$ is simple.\\
    If $\chi=\chi_U$ is irreducible, then $U$ is simple, so $U^*$ is simple by the first part. Proposition~\ref{prop:14.8} gives
    \[
        \chi_{U^*}=\overline{\chi_U}=\overline{\chi},
    \]
    and therefore $\overline{\chi}$ is again an irreducible character.
\end{solution}

If $F$ is any field and $U$ is any $FG$-module, then we define the character of $U$ to be an $F$-valued function on $U$ whose value at $g \in G$ is the trace of the $F$-endomorphism of $U$ induced by $g$, exactly as was done in this section in the case $F = \C$.

\begin{exercise}
\label{ex:14.7}
Suppose that $F$ has characteristic $p > 0$. Give an example to demonstrate that two $FG$-modules that have the same character need not be isomorphic.
\end{exercise}
\begin{solution}
    Let $G=\langle g\rangle$ be cyclic of order $p$, and let $F$ have characteristic $p$. Let $V=F^2$ with the trivial $FG$-module structure. Let $U=F^2$ with $g$ acting by the matrix
    \[
        J=
        \begin{pmatrix}
            1 & 1\\
            0 & 1
        \end{pmatrix}.
    \]
    Since $J=I+N$ with $N^2=0$, and since $F$ has characteristic $p$, we have
    \[
        J^p=(I+N)^p=I+N^p=I.
    \]
    Thus this defines an $FG$-module. For $0\leq k<p$,
    \[
        J^k=
        \begin{pmatrix}
            1 & k\\
            0 & 1
        \end{pmatrix},
    \]
    so $\operatorname{tr}(J^k)=2$. On the trivial module $V$, every element of $G$ acts as the $2\times2$ identity matrix, whose trace is also $2$. Therefore $\chi_U=\chi_V$.\\
    However, $U$ and $V$ are not isomorphic. The fixed-point space $V^G$ is all of $V$, so $\dim_F V^G=2$. On the other hand,
    \[
        U^G=\ker(J-I),
    \]
    which is one-dimensional. Since an isomorphism of $FG$-modules preserves the dimension of the fixed-point space, $U\not\cong V$.
\end{solution}

\begin{exercise}
\label{ex:14.8}
Suppose that $F$ has characteristic $p > 0$. Show that if $U$ is an $FG$-module then we have $\chi_U(g) = \chi_U(g_{p'})$ for any $g \in G$, where $g_{p'}$ is the $p'$-part of $g$ as defined in Exercise~\ref{ex:1.4}.
\end{exercise}
\begin{solution}
    Write
    \[
        g=xy=yx,
    \]
    where $x=g_p$ is the $p$-part of $g$ and $y=g_{p'}$ is the $p'$-part of $g$, as in Exercise~\ref{ex:1.4}. Let $X$ and $Y$ be the linear transformations of $U$ induced by $x$ and $y$, respectively. Since $x$ and $y$ commute, $X$ and $Y$ commute. Also $x$ has $p$-power order, so for some $a\geq0$ we have $X^{p^a}=I$. In characteristic $p$,
    \[
        (X-I)^{p^a}=X^{p^a}-I=0.
    \]
    Thus $N=X-I$ is nilpotent. Since $N$ commutes with $Y$, the product $NY$ is also nilpotent:
    \[
        (NY)^{p^a}=N^{p^a}Y^{p^a}=0.
    \]
    A nilpotent endomorphism has trace zero. Therefore
    \[
        \chi_U(g)
        =
        \operatorname{tr}(XY)
        =
        \operatorname{tr}((I+N)Y)
        =
        \operatorname{tr}(Y)+\operatorname{tr}(NY)
        =
        \operatorname{tr}(Y)
        =
        \chi_U(g_{p'}).
    \]
\end{solution}

\end{document}
